\section{Introduction}
\label{sec:intro}

The cultural or social meaning of a body posture is arbitrarily linked to its physical appearance.
Postures that look nearly identical can carry substantially different meanings.
A model that reads appearance alone cannot determine which meaning applies.
Recovering the meaning requires the measured configuration of the body.

When a vision--language model misses the subtle physical differences that separate such postures, it can endorse harmful form in physical therapy or exercise by confirming a misaligned joint as correct.
The same failure can misrepresent postures whose meaning matters to the communities practicing them.

Inferring meaning from appearance therefore requires explicit, measured body geometry expressed in language.
Human pose analysis is commonly formulated as keypoint localization~\cite{cao2019openpose,sun2019deep}, skeleton recovery, or action recognition.
These tasks provide perceptual representations but do not answer natural-language questions about configuration, such as which side is more flexed or how a geometric property should be described in words.
Such questions require a model to connect image evidence with explicit body measurements and linguistic reasoning.

Yoga is a particularly demanding setting for this problem.
Many asanas are distinguished by small differences in joint flexion, limb elevation, trunk orientation, and bilateral symmetry.
Yoga-82~\cite{verma2020yoga82} provides a large and visually diverse benchmark for fine-grained recognition across 82 pose categories.
Its supervision is limited to hierarchical category labels.
It does not provide reconstructed 3D bodies, interpretable geometric descriptors, or question--answer annotations that connect a specific image to its measured configuration.

Recent progress in monocular human mesh recovery and visual instruction tuning offers a path toward such supervision.
SAM 3D Body recovers articulated body structure under the Momentum Human Rig representation from a single image~\cite{yang2026sam3dbody}.
Models such as LLaVA learn to follow language instructions conditioned on visual input~\cite{liu2023visual}.
A direct combination, however, leaves two unresolved issues.
First, dense mesh outputs must be converted into compact and interpretable measurements.
Second, continuous geometric values must be integrated into language generation without reducing numeric prediction to unconstrained text-token generation.

We introduce \textbf{YogaAlign}, a benchmark that pairs Yoga-82 images with measured body geometry and question--answer annotations grounded in that geometry.
The geometry comes from SAM 3D Body reconstructions, summarized into a taxonomy of angular, symmetry, trunk, and support descriptors.
The annotations form two sets.
We produce pose-level domain-knowledge QA by evidence-conditioned synthesis from textual sources.
We generate geometry-grounded QA directly from the measurements of each image.

Building on YogaAlign, we train \textbf{YogaLLaVA V3}, a vision--language model that conditions on this geometry during generation.
Ordered \texttt{<GEO>} tokens carry the angular subset of the measurements.
The \texttt{<ANGLE>} mechanism separates the decision to place a numeric slot from continuous scalar regression and feeds the predicted value back into the autoregressive context.
In free-running evaluation, YogaLLaVA V3 produces the correct number of numeric measurements more often than the best baseline (86.1\% of samples versus 68.1\% for Qwen2.5-VL-7B) and achieves lower angular error on the shared common-success subset (a mean absolute error of $17.65^\circ$ versus $42.80^\circ$ for LLaVA-1.5-7B).

We make three contributions.
\begin{itemize}[leftmargin=1.2em,itemsep=2pt,topsep=3pt]
    \item We introduce \textbf{YogaAlign}, a geometry-grounded extension of Yoga-82 containing 3D-derived descriptors and 10,090 instruction-style question--answer pairs across 82 asanas.
    \item We develop a two-part annotation procedure that combines evidence-conditioned yoga-domain knowledge with image-specific geometric measurements, exact angular targets, rule-based validation, and traceable provenance.
    \item We build \textbf{YogaLLaVA V3} on YogaAlign, integrating metric-aware geometry tokens, continuous angle-slot regression over descriptors supplied through those tokens, and scalar feedback into a vision--language decoder, and evaluate both teacher-forced angle regression and free-running generation.
\end{itemize}
Our experiments cover the geometry-grounded set alone.
The domain-knowledge set is released as a training and evaluation resource, and we report no model performance on it.
